\section{Background: SNOMED~CT}
\label{sec:background}
\label{sec:snomed}

SNOMED~CT is the most comprehensive clinical terminology in use \cite{snomed2024starter}.
The release we use contains more than 420,000 concepts, covering disorders, procedures, substances, body structures and more, connected by 113 relationship types.
Concepts and relationships together form a DAG (Section~\ref{sec:method:problem}).

\paragraph{Hierarchy and attributes.}
Relationships are of two kinds.
An \isa{} relationship $u \xrightarrow{\isa} v$ states that $u$ is a subtype of $v$; a concept may have several parents, so the hierarchy is a polyhierarchy.
Every other relationship states the value of an attribute: $u \xrightarrow{R} v$, for $R$ such as \emph{finding site} or \emph{associated morphology}, means that the value of attribute $R$ for $u$ is $v$, not that $v$ is more general than $u$.
For instance, \emph{Neoplasm of bone of upper limb} is a neoplasm (its associated morphology) located in the bone of the upper limb (its finding site), and a subtype of \emph{Neoplasm of bone} and of \emph{Structural abnormality of upper limb bone} (Figure~\ref{fig:shared-structure}).
Our tokenizer treats the two kinds of edge differently (Section~\ref{sec:method:tokenizer}).

\paragraph{Shared meaning.}
Because meaning is carried by edges, two distinct concepts share part of their meaning when they reach the same concepts through the same relationships.
\emph{Fracture of upper limb} and \emph{Neoplasm of bone of upper limb} are different disorders, yet both are subtypes of \emph{Structural abnormality of upper limb bone} and both have \emph{Bone structure of upper limb} as their finding site (Figure~\ref{fig:shared-structure}).
This shared structure is what allows distinct concepts to share tokens: a rare concept can be described by the more frequent concepts it relates to.

\paragraph{Definition status.}
A \emph{fully defined} concept is equivalent to the conjunction of its relationships: any concept with the same relationships is the same concept.
For a \emph{primitive} concept, the relationships are necessary but not sufficient, and part of its meaning is stated only by its name \cite{snomed2024editorial}.
Replacing a primitive concept by the concepts it reaches therefore loses this unstated part, which bounds what any tokenization based on the graph can preserve (Sections~\ref{sec:method:score} and~\ref{sec:limitations}).

\begin{figure}[t]
\centering
\includegraphics[width=\linewidth]{paper/imgs/concept_relations.png}
\caption{Two SNOMED~CT disorders (blue) and their direct neighbours. Upward edges are \isa{} relationships to more general concepts; downward edges are attribute relationships, whose targets are values of the attribute rather than more general concepts. The two disorders differ, yet share a parent (\emph{Structural abnormality of upper limb bone}) and a finding site (\emph{Bone structure of upper limb}), shown in red.}
\label{fig:shared-structure}
\end{figure}
